%!TEX program = lualatex
%!TEX encoding = UTF-8 Unicode  
%!TEX root = chapter01.tex  
\documentclass[main.tex]{subfiles}
\begin{document} 
%----------------------------------------------------------------------------------------
%	CHAPTER 10
%----------------------------------------------------------------------------------------
\cleardoublepage
\loadgeometry{marginlayout}  
\pagestyle{scrheadings} % Print headers again  
\KOMAoptions{
	headwidth=textwithmarginpar,
	headsepline=true,
	headsepline= 0.5pt,
	footwidth=textwithmarginpar,
}   
\onehalfspacing 
\setcounter{chapter}{0}

%\begingroup
%% Set the style for section entries
%\etocsetstyle{section} 
%{\parindent -5pt \parskip 0pt}
%{\leftskip 0pt}
%{\makebox[.8cm]{\etocnumber\autodot}
%	\etocname\nobreak\leaders\hbox{\hbox to 1.5ex {\hss.\hss}}\hfill\nobreak
%	\etocpage\par}
%{}
%% Set the style for subsection entries
%\etocsetstyle{subsection}
%{\parindent -5pt \parskip 0pt}
%{\leftskip 0pt}
%{\makebox[.5cm]{}
%	\etocname\nobreak\leaders\hbox{\hbox to 1.5ex {\hss.\hss}}\hfill\nobreak
%	\etocpage\par}
%{}
%% Set the global style of the toc
%%\etocsettocstyle{}{}
%%\etocsettocstyle{\normalfont\sffamily\normalsize}{} 
%\etocsettocstyle{\usekomafont{section}\small}{}
%\etocsetnexttocdepth{\themargintocdepth}
%% Print the table of contents in the margin 
%\marginnote[#1]{}
%\endgroup


\chapter{Pourcentages et évolutions}

%---------------------------------------------------------------------------------------- 
%	 
%----------------------------------------------------------------------------------------
\begingroup
% Set the style for section entries
\etocsetstyle{section} 
{\parindent -5pt \parskip 0pt}
{\leftskip 0pt}
{\makebox[.8cm]{\etocnumber\autodot}
	\etocname\nobreak\leaders\hbox{\hbox to 1.5ex {\hss.\hss}}\hfill\nobreak
	\etocpage\par}
{}
% Set the style for subsection entries
\etocsetstyle{subsection}
{\parindent -5pt \parskip 0pt}
{\leftskip 0pt}
{\makebox[.5cm]{}
	\etocname\nobreak\leaders\hbox{\hbox to 1.5ex {\hss.\hss}}\hfill\nobreak
	\etocpage\par}
{}
% Set the global style of the toc
%\etocsettocstyle{}{}
%\etocsettocstyle{\normalfont\sffamily\normalsize}{} 
\etocsettocstyle{\usekomafont{section}\small}{}
\etocsetnexttocdepth{\themargintocdepth}
% Print the table of contents in the margin 
\localtableofcontents
\endgroup
%---------------------------------------------------------------------------------------- 
%	 
%----------------------------------------------------------------------------------------

\newpage
\section{Taux d'évolution et coefficient multiplicateur}

\begin{definition}[Le $P$\%] désigne $P$ centièmes$=\frac{P}{100}$. 
\end{definition}
\begin{definition}[$U$ de $V$] désigne $U\times V$.
\end{definition}
 \begin{proof} L'évolution de taux $TE$ appliquée à $X$ donne $X+TE\times X=(1+TE)X$.
	\end{proof}
\begin{definition}[taux d'évolution et coefficient multiplicateur] Une évolution de taux $TE$ correspond à un multiplication par $CM=1+TE$  
	\setlength{\abovedisplayskip}{3pt} 
	\setlength{\belowdisplayskip}{3pt} 
	$$
	\text{coefficient multiplicateur}  =  1 +\text{taux d'évolution}
	$$ 
\end{definition}
\begin{marginfigure} 
	\tikzstyle{noeud}=[rounded corners=3pt,fill=gray!35, text centered,text width=1.5cm]
	\begin{tikzpicture}
		\node[noeud](A)at(0,0){$V_I=x$}; 
		\node[noeud](C)at(3.5,0){$V_F=y$}; 
		\node(D)at(1.75,1.5){$\times CM =\frac{y}{x}$};
		\draw[->,>=latex](A)to[bend right=-30](C); 
		\node(E)at(2,-1.5){$\times (1+TE)$};
	\end{tikzpicture}
	\caption{Evolution, CM et TE}
	\label{fig:carreproduit}
\end{marginfigure}

 
\noindent
\begin{tabularx}{0.95\linewidth}{XXX}
\toprule
\textbf{Journalistes} & \textbf{TE} & \textbf{CM}$=1+TE$ \\ \toprule
augmentation de 12\% & 12\% =\numprint{0.12}& $CM=\numprint{1.12}$ \\ \hline
diminution de 12\% &  $-12\%=-\numprint{0.12}$ & \numprint{0.88} \\\hline
 pas de changement & $0$ & $1$ \\   \bottomrule 
\end{tabularx}

 \begin{proposition}Dans une évolution $V_I\mapsto V_F$ on a :
 	\[
 	V_F= CM\times V_I \qquad V_F = (1+TE) \times V_I
 	\]
 	
 	\[
 	CM= \frac{V_F}{V_I} \qquad TE = CM - 1 = \frac{V_F-V_I}{V_I}
 	\]
\end{proposition}


%----------------------------------------------------------------------------------------
%	EXERCICES
%----------------------------------------------------------------------------------------

\newpage 
\loadgeometry{widelayout}  
\pagestyle{scrheadings} % Print headers again  
\KOMAoptions{
	headwidth=textwithmarginpar,
	headsepline=true,
	headsepline= 0.5pt,
	footwidth=textwithmarginpar,
} 
\onehalfspacing
\def\espacecarre{\begin{tikzpicture}
		\node [rectangle,draw,dashed, minimum size=0.7cm]  at (2,-0.2) {\qquad\qquad};\end{tikzpicture}}
\subsection{Exercices : taux d'évolution et coefficient multiplicateur} 
\begin{exercise} Complétez.\vspace{-1em}
	\begin{spacing}{2}
		\begin{enumerate}[label=\alph*),leftmargin=*]
			\item  Une augmentation de $3\%$ est une évolution de taux $TE= $\dotfill.\\ Elle correspond à une multiplication par $CM=\ldots +\ldots\ldots =$\dotfill  
			\item Une diminution de $7\%$ est une évolution de taux $TE=-$\dotfill.\\
			 Elle correspond à une multiplication par $CM=\ldots +\ldots\ldots =$\dotfill  
			\item Multiplier par $CM=1.2$ correspond à une évolution de taux $TE=$\dotfill. \\
			C'est une (augmentation/diminution)  de $\ldots\ldots \%$.
			\item Multiplier par $CM=0.95$ correspond à une évolution de taux $TE=\ldots \ldots- \ldots\ldots=$\dotfill.\\
			 C'est une (augmentation/diminution)  de $\ldots\ldots  \%$.
			 \item Une évolution de \EUR{1.50} à \EUR{1.86} correspond à une multiplication par $CM=\frac{\espacecarre}{\espacecarre}=\espacecarre$. C'est une évolution de taux $TE=\espacecarre -1=\espacecarre\%$.
			 \item Une évolution de \EUR{40} à  EUR{24} correspond à une multiplication par $CM=\frac{\espacecarre}{\espacecarre}=\espacecarre$. C'est une évolution de taux $TE=\espacecarre-1=\espacecarre\%$.
			 \item Une évolution de \EUR{320} à \EUR{288} est de taux $TE=\frac{\espacecarre-\espacecarre}{\espacecarre\null}=\espacecarre=\espacecarre\%$.
			 \item Une évolution de \EUR{90} à \EUR{100} est de taux $TE=\frac{\espacecarre-\espacecarre}{\espacecarre\null}=\espacecarre=\espacecarre\%$.
		\end{enumerate} 
	\end{spacing}
\end{exercise}
\vspace{-2.5em}
\begin{eBox}
\noindent Les  exercices suivants illustrent des   formulations de   problèmes d'évolutions. 
\end{eBox}
\begin{exercise}\hfill 
	
	\noindent\parbox{0.70\linewidth}{
Prix initial est de \EUR{80}. Après augmentation le prix est de \EUR{125}. \\
Quel est le taux d'augmentation ? 
}\parbox{0.23\linewidth}{
\tikzstyle{noeud}=[rounded corners=3pt,fill=gray!35, text centered,text width=2.0cm]
\begin{tikzpicture}
	\node[noeud](A)at(0,0){$V_I=\qquad $}; 
	\node[noeud](C)at(3.5,0){$V_F=\qquad$}; 
%	\node(D)at(1.75,1.5){\scriptsize$\times CM =$};
%		\node(E)at(1.75,1.2){\scriptsize$\times \overbrace{(1+TE)}^{CM}=$};
		\node(E)at(1.75,1.0){\scriptsize$\times CM=$};
	\draw[->,>=latex](A)to[bend right=-30](C); 
	\node(F)at(1.2,-0.8){\scriptsize$TE=\qquad$};
	\node(F)at(1.2,-1.3){\scriptsize$\nearrow$ de}; 
\end{tikzpicture}
}	 
\end{exercise}
\begin{exercise}\hfill 
	
	\noindent\parbox{0.70\linewidth}{
	Le prix initial est de \EUR{16}. 
	
	Après réduction le prix est de \EUR{12.5}.
	
	Donner le taux de diminution   
	}\parbox{0.23\linewidth}{
\tikzstyle{noeud}=[rounded corners=3pt,fill=gray!35, text centered,text width=2.0cm]
		\begin{tikzpicture}
		\node[noeud](A)at(0,0){$V_I=\qquad $}; 
		\node[noeud](C)at(3.5,0){$V_F=\qquad$}; 
%		\node(E)at(1.75,1.2){\scriptsize$\times \overbrace{(1+TE)}^{CM}=$};
		\node(E)at(1.75,1.0){\scriptsize$\times CM=$};
		\draw[->,>=latex](A)to[bend right=-30](C); 
		\node(F)at(1.2,-0.8){\scriptsize$TE=\qquad$};
		\node(F)at(1.2,-1.3){\scriptsize$\nearrow$ de}; 
		\end{tikzpicture}
	}	 
\end{exercise}
 \begin{exercise}\hfill 
 	 
	\noindent\parbox{0.70\linewidth}{
Le montant de la redevance audiovisuel en France est passé de \EUR{114.49}  en 2001 à \EUR{123}  en 2011.\\
Quel est le taux d’évolution de cette taxe de 2001 à 2011 ?
}\parbox{0.23\linewidth}{
	\tikzstyle{noeud}=[rounded corners=3pt,fill=gray!35, text centered,text width=2.0cm]
	\begin{tikzpicture}
		\node[noeud](A)at(0,0){$V_I=\qquad $}; 
		\node[noeud](C)at(3.5,0){$V_F=\qquad$}; 
		%		\node(E)at(1.75,1.2){\scriptsize$\times \overbrace{(1+TE)}^{CM}=$};
		\node(E)at(1.75,1.0){\scriptsize$\times CM=$};
		\draw[->,>=latex](A)to[bend right=-30](C);
		\node(F)at(1.2,-0.8){\scriptsize$TE=\qquad$};
	\end{tikzpicture}
}
 \end{exercise}
\begin{exercise}\hfill 
	
	\noindent\parbox{0.70\linewidth}{
	Le prix initial de \EUR{60} subit une augmentation de 35\%. \\
	Sans calculer le montant de l'augmentation retrouver le prix final ? 
}\parbox{0.23\linewidth}{
\tikzstyle{noeud}=[rounded corners=3pt,fill=gray!35, text centered,text width=2.0cm]
\begin{tikzpicture}
	\node[noeud](A)at(0,0){$V_I=\qquad $}; 
	\node[noeud](C)at(3.5,0){$V_F=\qquad$}; 
%		\node(E)at(1.75,1.2){\scriptsize$\times \overbrace{(1+TE)}^{CM}=$};
	\node(E)at(1.75,1.0){\scriptsize$\times CM=$};
	\draw[->,>=latex](A)to[bend right=-30](C);
	\node(F)at(1.2,-0.8){\scriptsize$TE=\qquad$};
\end{tikzpicture}
}	
\end{exercise}
\begin{exercise}\hfill 
	
	\noindent\parbox{0.70\linewidth}{
Le prix initial de \EUR{35} subit une diminution de 60\%.\\
Sans calculer le montant de l'augmentation retrouver le prix final ?  
}\parbox{0.23\linewidth}{
\tikzstyle{noeud}=[rounded corners=3pt,fill=gray!35, text centered,text width=2.0cm]
\begin{tikzpicture}
\node[noeud](A)at(0,0){$V_I=\qquad $}; 
\node[noeud](C)at(3.5,0){$V_F=\qquad$}; 
%		\node(E)at(1.75,1.2){\scriptsize$\times \overbrace{(1+TE)}^{CM}=$};
\node(E)at(1.75,1.0){\scriptsize$\times CM=$};
\draw[->,>=latex](A)to[bend right=-30](C); 
\node(F)at(1.2,-0.8){\scriptsize$TE=\qquad$};
\end{tikzpicture}
}	 
\end{exercise}
\begin{exercise}\hfill 
	
	\noindent\parbox{0.70\linewidth}{
En appliquant une augmentation de 12.5\% du prix initial, le prix augmente de \EUR{15}. 
	\begin{enumerate}[label=\alph*), leftmargin=*]
		\item Donner une équation vérifiée par le prix initial. \dotfill 
		\item Quel est le prix initial.\dotfill
	\end{enumerate} 
 }\parbox{0.23\linewidth}{
\tikzstyle{noeud}=[rounded corners=3pt,fill=gray!35, text centered,text width=2.0cm]
\begin{tikzpicture}
\node[noeud](A)at(0,0){$V_I=\qquad $}; 
\node[noeud](C)at(3.5,0){$V_F=\qquad$}; 
%		\node(E)at(1.75,1.2){\scriptsize$\times \overbrace{(1+TE)}^{CM}=$};
\node(E)at(1.75,1.0){\scriptsize$\times CM=$};
\draw[->,>=latex](A)to[bend right=-30](C); 
\node(F)at(1.2,-0.8){\scriptsize$TE=\qquad$};
\end{tikzpicture}
}	 
\end{exercise}
\begin{exercise}\hfill 
	
	\noindent\parbox{0.70\linewidth}{
Après augmentation de 12.5\%, le prix final est de \EUR{45}.  
\begin{enumerate}[label=\alph*), leftmargin=*]
	\item Donner une équation vérifiée par le prix initial. \dotfill 
	\item Quel est le prix initial.\dotfill
\end{enumerate} 
 }\parbox{0.23\linewidth}{
\tikzstyle{noeud}=[rounded corners=3pt,fill=gray!35, text centered,text width=2.0cm]
\begin{tikzpicture}
\node[noeud](A)at(0,0){$V_I=\qquad $}; 
\node[noeud](C)at(3.5,0){$V_F=\qquad$}; 
%		\node(E)at(1.75,1.2){\scriptsize$\times \overbrace{(1+TE)}^{CM}=$};
\node(E)at(1.75,1.0){\scriptsize$\times CM=$};
\draw[->,>=latex](A)to[bend right=-30](C); 
\node(F)at(1.2,-0.8){\scriptsize$TE=\qquad$};
\end{tikzpicture}
}	 
\end{exercise} 

\begin{exercise} \hfill 
	
		\noindent\parbox{0.70\linewidth}{
	Après diminution de 12.5\%, le prix final est de \EUR{80.5}. 
	
	Quel était le prix initial ?
	 }\parbox{0.23\linewidth}{
	\tikzstyle{noeud}=[rounded corners=3pt,fill=gray!35, text centered,text width=2.0cm]
	\begin{tikzpicture}
		\node[noeud](A)at(0,0){$V_I=\qquad $}; 
		\node[noeud](C)at(3.5,0){$V_F=\qquad$}; 
		%		\node(E)at(1.75,1.2){\scriptsize$\times \overbrace{(1+TE)}^{CM}=$};
		\node(E)at(1.75,1.0){\scriptsize$\times CM=$};
		\draw[->,>=latex](A)to[bend right=-30](C); 
		\node(F)at(1.2,-0.8){\scriptsize$TE=\qquad$};
	\end{tikzpicture}
}
\end{exercise} 
\begin{exercise} \hfill 
	
 		\noindent\parbox{0.70\linewidth}{
	En appliquant une diminution de 15\%, le prix final est \EUR{100}. \\
Quel est le montant de la diminution ?  
 }\parbox{0.23\linewidth}{
 	\tikzstyle{noeud}=[rounded corners=3pt,fill=gray!35, text centered,text width=2.0cm]
 	\begin{tikzpicture}
 		\node[noeud](A)at(0,0){$V_I=\qquad $}; 
 		\node[noeud](C)at(3.5,0){$V_F=\qquad$}; 
 		%		\node(E)at(1.75,1.2){\scriptsize$\times \overbrace{(1+TE)}^{CM}=$};
 		\node(E)at(1.75,1.0){\scriptsize$\times CM=$};
 		\draw[->,>=latex](A)to[bend right=-30](C); 
 		\node(F)at(1.2,-0.8){\scriptsize$TE=\qquad$};
 	\end{tikzpicture}
 }
\end{exercise}  
\begin{exercise}\hfill 
 
\noindent\parbox{0.70\linewidth}{
Le prix d’un appareil ménager a augmenté de 15 \% en 2 ans. Il coûte maintenant \EUR{460}.\\
 Quel est montant de l'augmentation en deux ans ? 
}\parbox{0.23\linewidth}{
	\tikzstyle{noeud}=[rounded corners=3pt,fill=gray!35, text centered,text width=2.0cm]
	\begin{tikzpicture}
		\node[noeud](A)at(0,0){$V_I=\qquad $}; 
		\node[noeud](C)at(3.5,0){$V_F=\qquad$}; 
		%		\node(E)at(1.75,1.2){\scriptsize$\times \overbrace{(1+TE)}^{CM}=$};
		\node(E)at(1.75,1.0){\scriptsize$\times CM=$};
		\draw[->,>=latex](A)to[bend right=-30](C); 
		\node(F)at(1.2,-0.8){\scriptsize$TE=\qquad$};
	\end{tikzpicture}
}
\end{exercise}
\begin{exercise}\hfill 
	 
\noindent\parbox{0.70\linewidth}{
	Les prix des aliments ont diminué de 20\%. Un aliment coûte maintenant 240. Quel est le montant de la diminution? 
}\parbox{0.23\linewidth}{
	\tikzstyle{noeud}=[rounded corners=3pt,fill=gray!35, text centered,text width=2.0cm]
	\begin{tikzpicture}
		\node[noeud](A)at(0,0){$V_I=\qquad $}; 
		\node[noeud](C)at(3.5,0){$V_F=\qquad$}; 
		%		\node(E)at(1.75,1.2){\scriptsize$\times \overbrace{(1+TE)}^{CM}=$};
		\node(E)at(1.75,1.0){\scriptsize$\times CM=$};
		\draw[->,>=latex](A)to[bend right=-30](C); 
		\node(F)at(1.2,-0.8){\scriptsize$TE=\qquad$};
	\end{tikzpicture}
}
\end{exercise}
 
 \begin{eBox}
 	Dans un plan d'épargne, les intérêts sont dit \textbf{simples} lorsqu'ils sont calculés chaque année sur la base de la \textbf{somme placée au départ}. 
 \end{eBox}   
 
\begin{eBox}
	Dans un plan d'épargne, les intérêts sont dit \textbf{composés} lorsqu'ils sont calculés chaque année sur la base de la \textbf{somme totale  accumulée l'année précédente}. De manière générale, les gains ou pertes d'un placement sont exprimés en pourcentage par rapport à l'année écoulée. 
\end{eBox} 
\begin{example}\hfill 

\noindent
\parbox{0.48\linewidth}{
Un placement  \EUR{200} en intérêts \textbf{simples} de 3\%, rapporte $3\%\times 200$ chaque année.

Au terme de 5 ans, la somme épargnée est 
\setlength{\abovedisplayskip}{3pt}
\setlength{\belowdisplayskip}{3pt}
$$
200+ 5\times 0.03\times 200=200\times 1.15
$$
Le montant de départ est multiplié par $CM_\text{global}=1.15$, soit une augmentation globale $TE_\text{global}=15\%$.
}\hfill 
\parbox{0.48\linewidth}{
Un placement  \EUR{200} en intérêts \textbf{composés} de 3\%, augmente de $3\%$ chaque année. Le placement est multiplié par $1.03$ chaque année.

Au terme de 5 ans, la somme épargnée est 
\setlength{\abovedisplayskip}{3pt}
\setlength{\belowdisplayskip}{3pt}
$$
200\times 1.03^5\approx 200\times 1.159
$$
Le montant de départ est multiplié par $CM_\text{global}\approx 1.159$, soit une augmentation globale $TE_\text{global}\approx16\%$.
} 
\end{example}
\begin{exercise}  Complétez le tableau. 
	
	\noindent	
	\begin{tabularx}{1\columnwidth}{|c|c|c|X|X|} \hline
		Placement initial &  taux d'intêret  &  nombre d'années &  Calcul &   Montant final \\ \hline\hline  
		\rule{0pt}{1.8em}\EUR{2000} & $5\%$ composés & 6 & $2000\times 1.05^{\ldots}$ & \\ \hline
		\rule{0pt}{1.8em}\EUR{4000} & $2\%$ composés & 8 &  & \\ \hline
		\rule{0pt}{1.8em}\EUR{3500} & $6\%$ composés & 10 &   & \\ \hline
		\rule{0pt}{1.8em}\EUR{10000} & $1\%$ composés & 11 &   & \\ \hline
		\rule{0pt}{1.8em}  &  &  & $5000 \times 1.07^{14}$  & \\ \hline
		\rule{0pt}{1.8em}  &  &  & $600 \times 1.1^{8}$  & \\ \hline
		\rule{0pt}{1.8em}\EUR{4000}  &  $2\%$ simples& 5 &    & \\ \hline
	\end{tabularx} 
\end{exercise} 
\begin{exercise} Entourez les bonnes réponses.\hfill 
	
	\noindent \QCM[  Alterne, Stretch=1.0, Reponses=3, Titre,AlphT, Largeur=2.5cm]{%
		%		{$70$ est solution de l'équation $x^{300}=1$ d'inconnue $x$}&2,%
		{Un placement de \EUR{1000}  se déprécie de $7\%$ par an. Le montant après 5 ans est ...}&{$1000\times 0.07\times 5 $}&{$1000\times 0.07^5 $}&{$1000\times 0.93^5 $}&3,%
		%		{Un placement de \EUR{2000}  rapporte $4\%$ d'intérêts composés par an. Le montant après 5 ans est ...}&{$2000\times 1.04\times 5 $}&{$2000\times 1.04^5 $}&{$2000\times 4^5 $}&2,% 
		{Un placement de \EUR{100}  rapporte $8\%$ d'intérêts composés par an. Le montant après 10 ans est ...}&{$100\times 1.8\times 10 $}&{$100\times 1.08^{10} $}&{$100\times 0.8^{10} $}&2,%
		{Un placement de \EUR{100}  rapporte $3.5\%$ d'intérêts composés par an. Le montant après 2 ans est ...}&{$100\times 0.07 $}&{$100\times 1.07^{2} $}&{$107.1225$}&2,%
		{Un placement de \EUR{1000}  se déprécie de $9\%$ par an. Le montant après 10 ans est ...}&{$1000\times 0.91^{10} $}&{$1000\times 0.9^{10} $}&{$1000\times 1.09^{10}$}&1,%
		{Un placement de \EUR{5000}  se déprécie de $12\%$ par an. Le montant après 6 ans est ...}&{$5000\times 0.12^{6} $}&{$1000\times 0.88^{6} $}&{$\scriptsize 1000(1-6\times 0.12)$}&2,%
		%		{Un placement de \EUR{11000}  rapporte $4\%$ d'intérêts composés par an. Le total des intérêts après 4 ans est ...}&{$\approx$\EUR{12700}}&{$\approx$\EUR{1700}}&{$\approx$\EUR{1868}}&1,%
%		{Un placement de \EUR{1000}  rapporte  $10\%$ d'intérêts composés par an. Le placement initial est doublé  au bout de...}&{$6$ ans}&{$8$ ans}&{$10$ ans}&2,%
		%		{Un placement de \EUR{1000}  se déprécie de $5\%$ par an. Le placement  initial est diminuée de moitié au bout de...}&{$13$ ans}&{$14$ ans}&{$15$ ans}&2,%
	}
\end{exercise} 
\begin{exercise}\hfill \\
	 Quark souhaite investir \EUR{10000} sur 10 ans. Le placement $A$ rapporte $4\%$ taux d'intérêts simples. Le placement $B$ rapporte $3.5\%$ d'intérêts composés.\\
	 Quel placement est le plus avantageux ? Montrer les calculs.
\end{exercise}
\begin{exercise}\hfill \\
	Une banque propose un placement à intérêts composés. La première année est à 7\%, et les suivantes sont à 5\%. \\
Quelle est le \textbf{taux d'}augmentation globale d'un placement après 3 années ? 
\end{exercise}

\begin{exercise} \hfill \\
	 Un placement de \EUR{1000} sur 10 ans se déprécie de $4\%$ par an. Calculer le \textbf{taux de} diminution globale sur $10$ ans.
\end{exercise}




%----------------------------------------------------------------------------------------
 



\begin{exercise}[bilan]  	Compléter le tableau (Variation absolue = Prix final - Prix initial).  
	
%	\begin{table*}[h] 
	\noindent	\begin{tabularx}{1\columnwidth}{|X|c|c|c|c|c|} \hline
			\scriptsize Augmentation/diminution &\scriptsize taux d'évolution  &\scriptsize Coefficient Multiplicateur &\scriptsize Prix initial& \scriptsize Prix final &\scriptsize Variation absolue \\ \hline\hline  
			\rule[-0.6em]{0em}{1.9em} 	 &  +100\%   &&\EUR{98.40}&   &   \\ \hline
			\rule[-0.6em]{0em}{1.9em} 	 &  +160\%   &&\EUR{196.80} &  &   \\ \hline  
			%			\rule[-0.6em]{0em}{1.8em} \EUR{50}& +20\% & & & & \\ \hline
			\rule[-0.6em]{0em}{1.9em}  &+20\%  & & \EUR{60}&& \\ \hline
			\rule[-0.6em]{0em}{1.9em} & -20\%  & & \EUR{72}& & \\ \hline
			\rule[-0.6em]{0em}{1.9em}  & & & \EUR{72}& \EUR{54} & \\ \hline
			\rule[-0.6em]{0em}{1.9em}  Augmentation de 50\% &	&  & &\EUR{54} & \\ \hline
			\rule[-0.6em]{0em}{1.9em} Diminution de 50\%&	&     & &\EUR{54} & \\ \hline
			\rule[-0.6em]{0em}{1.9em} 	& +20\% &  & &\EUR{54} & \\ \hline
			\rule[-0.6em]{0em}{1.9em}  	& -20\% & & &\EUR{54} & \\ \hline
			\rule[-0.6em]{0em}{1.9em} 	   & & &\EUR{54}& & +\EUR{54}\\ \hline
			\rule[-0.6em]{0em}{1.9em} 	    & & $\times \numprint{0.7}$ &\EUR{40} &   &   \\ \hline 
			\rule[-0.6em]{0em}{1.9em} 	&   & & &\EUR{108} & $-$\EUR{27}\\ \hline
			\rule[-0.6em]{0em}{1.9em} 	&    &  $\times \numprint{1.3}$& & \EUR{91} &   \\ \hline 
			\rule[-0.6em]{0em}{1.9em} 	   & & &\EUR{96}&\EUR{108} &  \\ \hline
			\rule[-0.6em]{0em}{1.9em}  &+\numprint{1.25}\%   & &	\EUR{96}&  &  \\ \hline
			\rule[-0.6em]{0em}{1.9em} 	 &+\numprint{25}\% &  & &\EUR{98.40}  &  \\ \hline
			\rule[-0.6em]{0em}{1.9em} 	   & & $\times \numprint{1.007}$ &\EUR{130}   &   &   \\ \hline 
			\rule[-0.6em]{0em}{1.9em} 	    & & & \EUR{98.40}& &  $-$\EUR{19.68}\\ \hline 
			\rule[-0.6em]{0em}{1.9em} 	 &  +42\% & & \EUR{17} & &   \\ \hline  
		\end{tabularx}
%	\end{table*} 
\end{exercise} 


%----------------------------------------------------------------------------------------
%	cours
%----------------------------------------------------------------------------------------

\newpage 
\loadgeometry{marginlayout}  
\pagestyle{scrheadings} % Print headers again  
\KOMAoptions{
	headwidth=textwithmarginpar,
	headsepline=true,
	headsepline= 0.5pt,
	footwidth=textwithmarginpar,
} 
\onehalfspacing
\section{Evolutions successives}


\begin{theorem}[Évolutions successives]
	Plusieurs évolutions successives vont avoir le même effet qu'une seule dont le \textbf{CM global est le produit des CM  des évolutions intermédiaires qui la composent}.  
\end{theorem}  
\begin{theorem} [Évolutions successives de même TE] 
	Pour une succession de $n$ évolutions de  même TE, le taux d'évolution global est donné par :
	\setlength{\abovedisplayskip}{3pt}
	\setlength{\belowdisplayskip}{3pt}
	\begin{align*} 
		CM_{\text{global}}   &= (CM)^n \\
		1+TE_{\text{global}}  &= (1+TE )^n
	\end{align*} 
\end{theorem}
	\begin{marginfigure} 
		\tikzstyle{noeud}=[rounded corners=3pt,fill=gray!35, text centered,text width=1.2cm]
		\begin{tikzpicture}[scale=0.7]
			\node[noeud](A)at(0,0){$V_D=?$};
			\node[noeud](B)at(3,0){$V_I=?$};
			\node[noeud](C)at(6,0){$V_A=?$};
			\node(D)at(1.5,-1.5){$\times CM_1=\ldots$};
			\node(D1)at(1.5,-2.5){$TE_1=\ldots$};
			\node(D)at(4.5,-1.5){$\times CM_2=\ldots$};
			\node(D1)at(4.5,-2.5){$TE_2=\ldots$};
			\node(D)at(3,1.5){$\times CM_\text{global}=\ldots$};
			\node(D1)at(3,2.5){$TE_\text{global}=\ldots$};
			\draw[->,>=latex](A)to[bend right=-30](C);
			\draw[->,>=latex](B)to[bend right=60](C);
			\draw[->,>=latex](A)to[bend right=60](B);
		\end{tikzpicture}
	\caption{Pour deux évolutions successives   de CM $3$ et $7$, vont avoir le même effet qu'une seule  dont le CM est $3\times 7$.}
	\end{marginfigure} 
\begin{theorem}[Cas de 2   évolutions successives] 
	\setlength{\abovedisplayskip}{3pt}
	\setlength{\belowdisplayskip}{3pt}
\begin{align} 
	CM_{\text{global}}   &= CM_1 \times CM_2   \label{eq.successives.1}\\
	1+TE_{\text{global}}  &= (1+TE_1 ) \times (1+TE_2)  \label{eq.successives.2}
\end{align}
\end{theorem} 


\subsection{Taux d'évolution réciproque}
 
\begin{figure}[h]
	\tikzstyle{noeud}=[rounded corners=3pt,fill=gray!35, text centered,text width=1.2cm]
	\begin{tikzpicture}
		\node[noeud](A)at(0,0){$V_0$};
		\node[noeud, text width=2cm](B)at(3,0){$V_1$};
		\node[noeud](C)at(6,0){\small $V_2=V_0$};
		\node(D)at(1.5,-1.5){$\times CM_1=\frac{5}{4}$};
		\node(D)at(4.5,-1.5){$\times CM_2=\frac{4}{5}$};
		\node(D)at(3,1.5){$\times CM_\text{global}=1$};
		\draw[->,>=latex](A)to[bend right=-30](C);
		\draw[->,>=latex](B)to[bend right=60](C);
		\draw[->,>=latex](A)to[bend right=60](B);
		\node(D1)at(1.5,-2.5){$TE_1=+25\%$}; 
		\node(D1)at(4.5,-2.5){$TE_2=-20\%$}; 
		\node(D1)at(3,2.5){$TE_\text{global}=0$};
	\end{tikzpicture}
	\caption{L'évolution $100\mapsto 125$ et l'évolution $125\mapsto 100$ sont réciproques.}
\end{figure}
\begin{definition}[Évolution réciproque] de l'évolution $V_0\mapsto V_1$ est l'évolution $V_1\mapsto V_0$. 
	
	Les CM multiplicateurs sont inverses l'un de l'autre. 	Les formules \ref{eq.successives.1} et  \ref{eq.successives.2} donnent  : 
	\begin{align}  
		1 &  = CM \times CM_\text{reciproque} \\
		1 &  = (1+TE ) \times (1+TE_\text{reciproque}) 
	\end{align}
\end{definition}
\newpage 
\subsection{Approfondissement : Taux d'évolutions moyens}
\begin{definition}[CM et TE moyens] Deux évolutions successives de taux d'évolutions $TE_1$ et $TE_2$.
	
	On appelle \textbf{coefficient multiplicateur moyen} et \textbf{taux d'évolution moyen} les nombres : 
	\setlength{\abovedisplayskip}{3pt}
	\setlength{\belowdisplayskip}{3pt}
	\begin{align*} 
		CM_{\text{moyen}}^2  & = CM_1 \times CM_2 \\
		\big(1+TE_{\text{moyen}}\big)^2  &= (1+TE_1 ) \times (1+TE_2)
	\end{align*}  
	Donc $2$ évolutions successives de même taux $TE_{\text{moyen}}$ conduisent à une même évolution globale que les deux évolutions successives de $TE_1$ et $TE_2$.
\end{definition} 
\begin{proposition}[Cas de   3 évolutions successives] 
	\setlength{\abovedisplayskip}{3pt}
	\setlength{\belowdisplayskip}{3pt}
	\begin{align} 
		CM_{\text{global}}   &= CM_1 \times CM_2 \times CM_3 \label{eq.successives.3}\\
		1+TE_{\text{global}}  &= (1+TE_1 ) \times (1+TE_2) \times (1+TE_3)\label{eq.successives.4} \\
		\big(1+TE_{\text{moyen}}\big)^3  &= (1+TE_1 ) \times (1+TE_2)\times (1+TE_3)
	\end{align}
\end{proposition}  
\begin{example} Est-ce que trois augmentations de 12\%, puis 17\%, puis 10\% correspondent à une augmentation moyenne annuelle d'exactement 13\% ? 
\end{example}
%----------------------------------------------------------------------------------------
%	EXERCICES
%----------------------------------------------------------------------------------------

\newpage 
\loadgeometry{widelayout}  
\pagestyle{scrheadings} % Print headers again  
\KOMAoptions{
	headwidth=textwithmarginpar,
	headsepline=true,
	headsepline= 0.5pt,
	footwidth=textwithmarginpar,
} 
\onehalfspacing
\subsection{Exercices évolutions successives et réciproques} 
\begin{exercise} Déterminez le taux d'évolution global associé à la succession d'évolutions :
	\begin{enumerate}[label=\arabic*), leftmargin=*]
		\item augmentation de $8\%$, suivie d'une augmentation de $10\%$ \\
\noindent
\parbox{0.48\linewidth}{
 
$\begin{WithArrows}
CM_{\text{global}} &= CM_1\times CM_2 \rule{0pt}{1.9em}\\
1+TE_{\text{global}}& = (1+TE_1)(1+TE_2) \rule{0pt}{1.9em}\\
 1+TE_{\text{global}} &= \ldots\ldots\ldots  \times \ldots\ldots\ldots \rule{0pt}{1.8em} =\ldots\ldots\rule{0pt}{1.8em}\\
TE_{\text{global}} &=   \ldots\ldots\ldots-1 = \rule{0pt}{1.9em} 
\end{WithArrows}$  
}\hfill 
\parbox{0.48\linewidth}{ 
		\tikzstyle{noeud}=[rounded corners=3pt,fill=gray!35, text centered,text width=1.2cm]
\begin{tikzpicture}
	\node[noeud](A)at(0,0){$V_D=?$};
	\node[noeud](B)at(3,0){$V_I=?$};
	\node[noeud](C)at(6,0){$V_A=?$};
%	\node(D)at(1.5,-2.5){$\times CM_1=\ldots\ldots$};
	\node(D1)at(1.5,-1.5){$TE_1=\ldots\ldots$};
%	\node(D)at(4.5,-2.5){$\times CM_2=\ldots\ldots$};
	\node(D1)at(4.5,-1.5){$TE_2=\ldots\ldots$};
%	\node(D)at(3,2.5){$\times CM_\text{global}=\ldots\ldots$};
	\node(D1)at(3,1.5){$TE_\text{global}=\ldots\ldots$};
	\draw[->,>=latex](A)to[bend right=-30](C);
	\draw[->,>=latex](B)to[bend right=60](C);
	\draw[->,>=latex](A)to[bend right=60](B);
\end{tikzpicture}\\	
} \medskip\\
L'évolution globale correspond à une \dotfill  de $\ldots\ldots\%$.
 

	\item augmentation de 10\%  suivie d'une augmentation de 5\%  \\
\noindent
\parbox{0.58\linewidth}{ 
	$\begin{WithArrows} 
		1+TE_{\text{global}}& = (1+TE_1)(1+TE_2) \rule{0pt}{1.9em}\\
		1+TE_{\text{global}} &= \ldots\ldots\ldots  \times \ldots\ldots\ldots  =\ldots\ldots\rule{0pt}{1.9em} \\
		TE_{\text{global}} &=   \ldots\ldots\ldots-1 = \rule{0pt}{1.9em} 
	\end{WithArrows}$  
}\hfill 
\parbox{0.40\linewidth}{ 
	\tikzstyle{noeud}=[rounded corners=3pt,fill=gray!35, text centered,text width=1.2cm]
	\begin{tikzpicture}
		\node[noeud](A)at(0,0){$V_D$};
		\node[noeud](B)at(3,0){$V_I$};
		\node[noeud](C)at(6,0){$V_A$};
%		\node(D)at(1.5,-2.5){$\times CM_1=\ldots\ldots$};
		\node(D1)at(1.5,-1.5){$TE_1=\ldots\ldots$};
%		\node(D)at(4.5,-2.5){$\times CM_2=\ldots\ldots$};
		\node(D1)at(4.5,-1.5){$TE_2=\ldots\ldots$};
%		\node(D)at(3,2.5){$\times CM_\text{global}=\ldots\ldots$};
		\node(D1)at(3,1.5){$TE_\text{global}=\ldots\ldots$};
		\draw[->,>=latex](A)to[bend right=-30](C);
		\draw[->,>=latex](B)to[bend right=60](C);
		\draw[->,>=latex](A)to[bend right=60](B);
	\end{tikzpicture}\\	
}\\
L'évolution globale correspond à une \dotfill de $\ldots\ldots\%$
	\item diminution de 30\%  suivie d'une diminution de 10\%  \\
\noindent
\parbox{0.58\linewidth}{ 
	$\begin{WithArrows} 
		1+TE_{\text{global}}& =   \rule{0pt}{1.9em}\\
		1+TE_{\text{global}} &= \ldots\ldots\ldots  \times \ldots\ldots\ldots  =\ldots\ldots\rule{0pt}{1.9em} \\
		TE_{\text{global}} &=   \ldots\ldots\ldots-1 = \rule{0pt}{1.9em} 
	\end{WithArrows}$  
}\hfill 
\parbox{0.40\linewidth}{ 
	\tikzstyle{noeud}=[rounded corners=3pt,fill=gray!35, text centered,text width=1.2cm]
	\begin{tikzpicture}
		\node[noeud](A)at(0,0){$V_D$};
		\node[noeud](B)at(3,0){$V_I$};
		\node[noeud](C)at(6,0){$V_A$};
		%		\node(D)at(1.5,-2.5){$\times CM_1=\ldots\ldots$};
		\node(D1)at(1.5,-1.5){$TE_1=\ldots\ldots$};
		%		\node(D)at(4.5,-2.5){$\times CM_2=\ldots\ldots$};
		\node(D1)at(4.5,-1.5){$TE_2=\ldots\ldots$};
		%		\node(D)at(3,2.5){$\times CM_\text{global}=\ldots\ldots$};
		\node(D1)at(3,1.5){$TE_\text{global}=\ldots\ldots$};
		\draw[->,>=latex](A)to[bend right=-30](C);
		\draw[->,>=latex](B)to[bend right=60](C);
		\draw[->,>=latex](A)to[bend right=60](B);
	\end{tikzpicture}\\	
}\\
L'évolution globale correspond à une \dotfill  de $\ldots\ldots\%$
	\item diminution de 40\%  suivie d'une diminution de 10\% \\
\noindent
\parbox{0.58\linewidth}{ 
	$\begin{WithArrows} 
		1+TE_{\text{global}}& =  \rule{0pt}{1.9em}\\
		1+TE_{\text{global}} &= \ldots\ldots\ldots  \times \ldots\ldots\ldots  =\ldots\ldots\rule{0pt}{1.9em} \\
		TE_{\text{global}} &=   \ldots\ldots\ldots-1 = \rule{0pt}{1.9em} 
	\end{WithArrows}$  
}\hfill 
\parbox{0.40\linewidth}{ 
	\tikzstyle{noeud}=[rounded corners=3pt,fill=gray!35, text centered,text width=1.2cm]
	\begin{tikzpicture}
		\node[noeud](A)at(0,0){$V_D$};
		\node[noeud](B)at(3,0){$V_I$};
		\node[noeud](C)at(6,0){$V_A$};
		%		\node(D)at(1.5,-2.5){$\times CM_1=\ldots\ldots$};
		\node(D1)at(1.5,-1.5){$TE_1=\ldots\ldots$};
		%		\node(D)at(4.5,-2.5){$\times CM_2=\ldots\ldots$};
		\node(D1)at(4.5,-1.5){$TE_2=\ldots\ldots$};
		%		\node(D)at(3,2.5){$\times CM_\text{global}=\ldots\ldots$};
		\node(D1)at(3,1.5){$TE_\text{global}=\ldots\ldots$};
		\draw[->,>=latex](A)to[bend right=-30](C);
		\draw[->,>=latex](B)to[bend right=60](C);
		\draw[->,>=latex](A)to[bend right=60](B);
	\end{tikzpicture}\\	
}\\
L'évolution globale correspond à une \dotfill  de $\ldots\ldots\%$

\newpage
	\item diminution de 25\%  suivie d'une diminution de 20\%  \\
\noindent
\parbox{0.58\linewidth}{ 
	$\begin{WithArrows} 
		 & =   \rule{0pt}{1.9em}\\
	  &= \ldots\ldots\ldots  \times \ldots\ldots\ldots  =\ldots\ldots\rule{0pt}{1.9em} \\
		TE_{\text{global}} &=   \ldots\ldots\ldots-1 = \rule{0pt}{1.9em} 
	\end{WithArrows}$  
}\hfill 
\parbox{0.40\linewidth}{ 
	\tikzstyle{noeud}=[rounded corners=3pt,fill=gray!35, text centered,text width=1.2cm]
	\begin{tikzpicture}[scale=0.80]
		\node[noeud](A)at(0,0){$V_D$};
		\node[noeud](B)at(3,0){$V_I$};
		\node[noeud](C)at(6,0){$V_A$};
		%		\node(D)at(1.5,-2.5){$\times CM_1=\ldots\ldots$};
		\node(D1)at(1.5,-1.5){$TE_1=\ldots\ldots$};
		%		\node(D)at(4.5,-2.5){$\times CM_2=\ldots\ldots$};
		\node(D1)at(4.5,-1.5){$TE_2=\ldots\ldots$};
		%		\node(D)at(3,2.5){$\times CM_\text{global}=\ldots\ldots$};
		\node(D1)at(3,1.5){$TE_\text{global}=\ldots\ldots$};
		\draw[->,>=latex](A)to[bend right=-30](C);
		\draw[->,>=latex](B)to[bend right=60](C);
		\draw[->,>=latex](A)to[bend right=60](B);
	\end{tikzpicture}\\	
}\\
L'évolution globale correspond à une \dotfill  de $\ldots\ldots\%$
	\item augmentation de 25\%  suivie d'une diminution de 20\% \\
\noindent
\parbox{0.58\linewidth}{ 
	$\begin{WithArrows} 
		 & =   \rule{0pt}{1.9em}\\
	  &= \rule{0pt}{1.9em} \\
		TE_{\text{global}} &=     \rule{0pt}{1.9em} 
	\end{WithArrows}$  
}\hfill 
\parbox{0.40\linewidth}{ 
	\tikzstyle{noeud}=[rounded corners=3pt,fill=gray!35, text centered,text width=1.2cm]
	\begin{tikzpicture}[scale=0.80]
		\node[noeud](A)at(0,0){$V_D$};
		\node[noeud](B)at(3,0){$V_I$};
		\node[noeud](C)at(6,0){$V_A$};
		%		\node(D)at(1.5,-2.5){$\times CM_1=\ldots\ldots$};
		\node(D1)at(1.5,-1.5){$TE_1=\ldots\ldots$};
		%		\node(D)at(4.5,-2.5){$\times CM_2=\ldots\ldots$};
		\node(D1)at(4.5,-1.5){$TE_2=\ldots\ldots$};
		%		\node(D)at(3,2.5){$\times CM_\text{global}=\ldots\ldots$};
		\node(D1)at(3,1.5){$TE_\text{global}=\ldots\ldots$};
		\draw[->,>=latex](A)to[bend right=-30](C);
		\draw[->,>=latex](B)to[bend right=60](C);
		\draw[->,>=latex](A)to[bend right=60](B);
	\end{tikzpicture}\\	
}\\
L'évolution globale correspond à une \dotfill  de $\ldots\ldots\%$
	\item augmentation de 25\%  suivie d'une diminution de 25\% \\
\noindent
\parbox{0.58\linewidth}{ 
	$\begin{WithArrows} 
		 & =   \rule{0pt}{1.9em}\\
	  &=  \rule{0pt}{1.9em} \\
		TE_{\text{global}} &=     \rule{0pt}{1.9em} 
	\end{WithArrows}$  
}\hfill 
\parbox{0.40\linewidth}{ 
	\tikzstyle{noeud}=[rounded corners=3pt,fill=gray!35, text centered,text width=1.2cm]
	\begin{tikzpicture}[scale=0.80]
		\node[noeud](A)at(0,0){$V_D$};
		\node[noeud](B)at(3,0){$V_I$};
		\node[noeud](C)at(6,0){$V_A$};
		%		\node(D)at(1.5,-2.5){$\times CM_1=\ldots\ldots$};
		\node(D1)at(1.5,-1.5){$TE_1=\ldots\ldots$};
		%		\node(D)at(4.5,-2.5){$\times CM_2=\ldots\ldots$};
		\node(D1)at(4.5,-1.5){$TE_2=\ldots\ldots$};
		%		\node(D)at(3,2.5){$\times CM_\text{global}=\ldots\ldots$};
		\node(D1)at(3,1.5){$TE_\text{global}=\ldots\ldots$};
		\draw[->,>=latex](A)to[bend right=-30](C);
		\draw[->,>=latex](B)to[bend right=60](C);
		\draw[->,>=latex](A)to[bend right=60](B);
	\end{tikzpicture}\\	
}\\
L'évolution globale correspond à une \dotfill  de $\ldots\ldots\%$
	\item diminution de 20\%  suivie d'une augmentation de 25\%  \\
\noindent
\parbox{0.58\linewidth}{ 
	$\begin{WithArrows} 
	& =   \rule{0pt}{1.9em}\\
	&=  \rule{0pt}{1.9em} \\
	TE_{\text{global}} &=     \rule{0pt}{1.8em} 
\end{WithArrows}$  
}\hfill 
\parbox{0.40\linewidth}{ 
	\tikzstyle{noeud}=[rounded corners=3pt,fill=gray!35, text centered,text width=1.2cm]
	\begin{tikzpicture}[scale=0.80]
		\node[noeud](A)at(0,0){$V_D$};
		\node[noeud](B)at(3,0){$V_I$};
		\node[noeud](C)at(6,0){$V_A$};
		%		\node(D)at(1.5,-2.5){$\times CM_1=\ldots\ldots$};
		\node(D1)at(1.5,-1.5){$TE_1=\ldots\ldots$};
		%		\node(D)at(4.5,-2.5){$\times CM_2=\ldots\ldots$};
		\node(D1)at(4.5,-1.5){$TE_2=\ldots\ldots$};
		%		\node(D)at(3,2.5){$\times CM_\text{global}=\ldots\ldots$};
		\node(D1)at(3,1.5){$TE_\text{global}=\ldots\ldots$};
		\draw[->,>=latex](A)to[bend right=-30](C);
		\draw[->,>=latex](B)to[bend right=60](C);
		\draw[->,>=latex](A)to[bend right=60](B);
	\end{tikzpicture}\\	
}\\
L'évolution globale correspond à une \dotfill  de $\ldots\ldots\%$
	\item augmentation de 22\%  suivie d'une diminution de 15\% \\
\noindent
\parbox{0.58\linewidth}{ 
	$\begin{WithArrows} 
	& =   \rule{0pt}{1.9em}\\
	&=  \rule{0pt}{1.9em} \\
	TE_{\text{global}} &=     \rule{0pt}{1.9em} 
\end{WithArrows}$   
}\hfill 
\parbox{0.40\linewidth}{ 
	\tikzstyle{noeud}=[rounded corners=3pt,fill=gray!35, text centered,text width=1.2cm]
	\begin{tikzpicture}[scale=0.80]
		\node[noeud](A)at(0,0){$V_D$};
		\node[noeud](B)at(3,0){$V_I$};
		\node[noeud](C)at(6,0){$V_A$};
		%		\node(D)at(1.5,-2.5){$\times CM_1=\ldots\ldots$};
		\node(D1)at(1.5,-1.5){$TE_1=\ldots\ldots$};
		%		\node(D)at(4.5,-2.5){$\times CM_2=\ldots\ldots$};
		\node(D1)at(4.5,-1.5){$TE_2=\ldots\ldots$};
		%		\node(D)at(3,2.5){$\times CM_\text{global}=\ldots\ldots$};
		\node(D1)at(3,1.5){$TE_\text{global}=\ldots\ldots$};
		\draw[->,>=latex](A)to[bend right=-30](C);
		\draw[->,>=latex](B)to[bend right=60](C);
		\draw[->,>=latex](A)to[bend right=60](B);
	\end{tikzpicture}\\	
}\\
L'évolution globale correspond à une \dotfill  de $\ldots\ldots\%$
\end{enumerate} 
%\scriptsize augmentation de 200\% 
%& \scriptsize augmentation de 300\%   
	\end{exercise} 
\newpage 
\begin{exercise}[3 évolutions successives] 	Dans chaque cas, calculer le TE global correspondants aux évolutions successives ci-dessous :
\begin{enumerate}[label={\arabic*)},leftmargin=*] 
	\item une diminution de 30\% suivie d'une augmentation de 20\% suivie d'une diminution de 10\% 
	\item  une augmentation de 12\% suivie de deux baisses successives de 5\%
	\item trois augmentations successives de 10\%.
	\item trois diminutions successives de 5\%. 
	\end{enumerate} 
\end{exercise}

\begin{eBox}  Deux évolutions successives de taux d'évolutions $TE_1$ et $TE_2$.\\
	 On appelle   \textbf{taux d'évolution moyen} le nombre vérifiant : 
	\setlength{\abovedisplayskip}{3pt} 
	\setlength{\belowdisplayskip}{3pt} 
	\begin{gather*} 
	\big(1+TE_{\text{moyen}}\big)^2   = (1+TE_1 )   (1+TE_2)\\
	1+TE_{\text{moyen}}=\sqrt{ (1+TE_1 )   (1+TE_2)}
	\end{gather*}
	Donc $2$ évolutions successives de même taux $TE_{\text{moyen}}$ conduisent à une même évolution globale que les deux évolutions successives de $TE_1$ et $TE_2$.
%	\\
%Le concept se généralise à 3 évolutions de taux $TE_1$, $TE_2$ et $TE_3$ à l'aide de la racine cubique :
%	\setlength{\abovedisplayskip}{3pt} 
%	\setlength{\belowdisplayskip}{3pt} 
%	\begin{gather*} 
%		\big(1+TE_{\text{moyen}}\big)^3   = (1+TE_1 )   (1+TE_2) (1+TE_3)\\
%		1+TE_{\text{moyen}}=\sqrt[3]{ (1+TE_1 )   (1+TE_2)  (1+TE_3)}
%		\end{gather*}
\end{eBox}  

\begin{exercise} \hfill 
	
	\noindent		Le nombre de nouvelles inscriptions Netflix à augmenté de 21\% durant le moi d'octobre, puis de 36\% durant le mois de novembre.
	Calculer le taux évolution mensuel moyen.
\end{exercise}  

\begin{exercise}\hfill 
	
	\noindent	 Le prix d'un article  augmente de $20\%$ la première année, et diminue de $4\%$ l'année suivante. Calculer le taux d'évolution annuel moyen.
\end{exercise}  
\begin{exercise} \hfill 
	
	\noindent Le montant d'un placement passe de \EUR{2000} en 2020 à   \EUR{2142.45} en 2022. Calculer le taux d'évolution annuel moyen.
\end{exercise} 
\begin{eBox} 
Le concept de taux moyen se généralise à 3 évolutions de taux $TE_1$, $TE_2$ et $TE_3$ à l'aide de la racine cubique :
\setlength{\abovedisplayskip}{3pt} 
\setlength{\belowdisplayskip}{3pt} 
\begin{gather*} 
	\big(1+TE_{\text{moyen}}\big)^3   = (1+TE_1 )   (1+TE_2) (1+TE_3)\\
	1+TE_{\text{moyen}}=\sqrt[3]{ (1+TE_1 )   (1+TE_2)  (1+TE_3)}
\end{gather*}
\end{eBox} 
\begin{exercise} \hfill 
	
	\noindent La taxe d'importations de patates-douces égyptiennes a augmenté de 9,75\% entre 2020 et 2022. Calculer le taux d'évolution annuel moyen de cette taxe.  
\end{exercise} 
\begin{exercise}\hfill 
	
	En Juin 2017, Farid achète une voiture pour \EUR{12000}. En Juin 2020, sa voiture est évaluée à \EUR{8600}.
	\begin{enumerate}[label={\alph*)},leftmargin=*] 
		\item Calculer le taux d'évolution annuel moyen.
		\item Si la voiture se déprécie à ce même taux. Quel est la valeur de la voiture en Juin 2022 ?
	\end{enumerate} 
\end{exercise}

\begin{exercise}\hfill 
	
	\noindent	Une moto achetée \EUR{2300} est revendue 3 ans plus tard à \EUR{1300}. Calculer le taux de dépréciation annuel moyen.
\end{exercise}

%-------------------------------------------------------------------------------
\newpage 
\begin{example}[Taux d'évolution réciproque] Pour chacune des évolutions suivantes, donner le taux d'évolution réciproque. Arrondir à $10^{-4}$ près si nécsssaire. 
\begin{enumerate}[label=\arabic*),leftmargin=*]
	\item augmentation de 25\% \\
\noindent
\parbox{0.58\linewidth}{ 
	$\begin{WithArrows} 
	1	& =  (1+TE)(1+TE_{\text{reciproque}}) \rule{0pt}{1.9em}\\
		&=  \rule{0pt}{1.9em} \\
		&=  \rule{0pt}{1.9em} \\
		TE_{\text{reciproque}} &=     \rule{0pt}{1.9em} 
	\end{WithArrows}$   
}\hfill 
\parbox{0.40\linewidth}{ 
	\tikzstyle{noeud}=[rounded corners=3pt,fill=gray!35, text centered,text width=1.2cm]
	\begin{tikzpicture}[scale=0.80]
		\node[noeud](A)at(0,0){$V_D$};
		\node[noeud](B)at(3,0){$V_I$};
		\node[noeud](C)at(6,0){$V_A$};
		%		\node(D)at(1.5,-2.5){$\times CM_1=\ldots\ldots$};
		\node(D1)at(1.5,-1.5){$TE=\ldots\ldots$};
		%		\node(D)at(4.5,-2.5){$\times CM_2=\ldots\ldots$};
		\node(D1)at(4.5,-1.5){$TE_{\text{reciproque}}=\ldots\ldots$};
				\node(D)at(3,1.5){$\times CM_\text{global}=1$};
%		\node(D1)at(3,2.5){$TE_\text{global}=\ldots\ldots$};
		\draw[->,>=latex](A)to[bend right=-30](C);
		\draw[->,>=latex](B)to[bend right=60](C);
		\draw[->,>=latex](A)to[bend right=60](B);
	\end{tikzpicture}\\	
}\\
L'évolution réciproque d'une  augmentation de $25\%$ est une \dotfill de $\ldots\ldots\%$.
	\item diminution de 25\% \\
\noindent
\parbox{0.58\linewidth}{ 
	$\begin{WithArrows} 
		1	& =  (1+TE)(1+TE_{\text{reciproque}}) \rule{0pt}{1.9em}\\
		&=  \rule{0pt}{1.9em} \\
		&=  \rule{0pt}{1.9em} \\
		TE_{\text{reciproque}} &=     \rule{0pt}{1.9em} 
	\end{WithArrows}$   
}\hfill 
\parbox{0.40\linewidth}{ 
	\tikzstyle{noeud}=[rounded corners=3pt,fill=gray!35, text centered,text width=1.2cm]
	\begin{tikzpicture}[scale=0.80]
		\node[noeud](A)at(0,0){$V_D$};
		\node[noeud](B)at(3,0){$V_I$};
		\node[noeud](C)at(6,0){$V_A$};
		%		\node(D)at(1.5,-2.5){$\times CM_1=\ldots\ldots$};
		\node(D1)at(1.5,-1.5){$TE=\ldots\ldots$};
		%		\node(D)at(4.5,-2.5){$\times CM_2=\ldots\ldots$};
		\node(D1)at(4.5,-1.5){$TE_{\text{reciproque}}=\ldots\ldots$};
		\node(D)at(3,1.5){$\times CM_\text{global}=1$};
		%		\node(D1)at(3,2.5){$TE_\text{global}=\ldots\ldots$};
		\draw[->,>=latex](A)to[bend right=-30](C);
		\draw[->,>=latex](B)to[bend right=60](C);
		\draw[->,>=latex](A)to[bend right=60](B);
	\end{tikzpicture}\\	
}\\
L'évolution réciproque d'une  diminution de $25\%$ est une \dotfill de $\ldots\ldots\%$.
\end{enumerate} 
\end{example}
\begin{exercise} Même consignes\vspace{-1em}
	\begin{multicols}{3} 
		\begin{enumerate}[label={\arabic*)},leftmargin=*] 
			\item augmentation de 50\%
			\item diminution de 20\%
			\item diminution de 90\%
			%			\item augmentation de 200\%
			\item augmentation de 300\%
			\item augmentation de $10$\%
			\item diminution de $10$\%
%			\item augmentation de $33$\%
%			\item augmentation de $13$\%
%			\item diminution de $27,2$\%
		\end{enumerate}
	\end{multicols}
\end{exercise}
\begin{exercise} \hfill 
	
	\noindent	Le prix TTC est 20\% du prix HT. Quelle évolution appliquer au prix TTC pour obtenir le prix HT ?
\end{exercise}


\begin{exercise}[Indice de base 100]\hfill 
	
	Le tableau ci-dessous donne le chiffre d'affaire annuel d'une entreprise pour les années comprises entre 2015 et 2021. 
	
	\noindent \begin{tabularx}{0.95\linewidth}{|c|X|X|X|X|X|X|X|}\hline 
		\textbf{Année} &2015 &2016 &2017 &2018 & 2019 & 2020 &2021 \\ \hline 
		\textbf{Chiffre d'affaire en miliers d'euros} & 134 & 138 & 138.3 & 135.6 & 133.2 & 138.2 &140.4 \\ \hline 
		\textbf{Indice (base 100)} & 97.1 & 100 & & &&&
		\\ \hline 
	\end{tabularx}
	\begin{enumerate}[label=\arabic*), leftmargin=*]
		\item Complétez la ligne des indices sachant qu'elle est proportionnellle à celle des chiffre d'affaire.
		\item Sans aucun calculs supplémentaires donner le taux d'augmentation en \% qui permet de passer du chiffre d'affaire de 2016 à celui de 2021.
		\item Même question avec le taux de diminution du chiffre d'affaire de 2016 à 2019.
	\end{enumerate}
\end{exercise}


%----------------------------------------------------------------------------------------
%	cours
%----------------------------------------------------------------------------------------
%
%\newpage 
%\loadgeometry{marginlayout}  
%\pagestyle{scrheadings} % Print headers again  
%\KOMAoptions{
%	headwidth=textwithmarginpar,
%	headsepline=true,
%	headsepline= 0.5pt,
%	footwidth=textwithmarginpar,
%} 
\onehalfspacing
%\section{Taux d'évolution réciproque}
% 
%
%
%
%
%
%\begin{figure}[h]
%	\tikzstyle{noeud}=[rounded corners=3pt,fill=gray!35, text centered,text width=1.2cm]
%	\begin{tikzpicture}
%		\node[noeud](A)at(0,0){$V_D=3$};
%		\node[noeud, text width=2cm](B)at(3,0){$V_A=15$};
%		\node[noeud](C)at(6,0){$V_D=3$};
%		\node(D)at(1.5,-1.5){$\times CM_1=5$};
%		\node(D)at(4.5,-1.5){$\times CM_2=\frac{1}{5}$};
%		\node(D)at(3,1.5){$\times CM_\text{global}=1$};
%		\draw[->,>=latex](A)to[bend right=-30](C);
%		\draw[->,>=latex](B)to[bend right=60](C);
%		\draw[->,>=latex](A)to[bend right=60](B);
%	\end{tikzpicture}
%	\caption{L'évolution $3\mapsto 15$ et l'évolution $15\mapsto 3$ sont réciproques.}
%\end{figure}
%\begin{definition}[Évolution réciproque] de l'évolution $V_D\mapsto V_A$ est l'évolution $V_A\mapsto V_D$. 
%	
%	Les CM multiplicateurs sont inverses l'un de l'autre. 	Les formules \ref{eq.successives.1} et  \ref{eq.successives.2} donnent  : 
%	\begin{align}  
%		1 &  = CM \times CM_\text{reciproque} \\
%		1 &  = (1+TE ) \times (1+TE_\text{reciproque}) 
%	\end{align}
%\end{definition}

%----------------------------------------------------------------------------------------
%	cours
%----------------------------------------------------------------------------------------

\newpage 
\loadgeometry{widelayout}  
\pagestyle{scrheadings} % Print headers again  
\KOMAoptions{
	headwidth=textwithmarginpar,
	headsepline=true,
	headsepline= 0.5pt,
	footwidth=textwithmarginpar,
} 
\onehalfspacing
\section{Club de Maths : Problèmes}
 

\begin{exercise}\hfill
	
	\noindent Une augmentation de 5\% suivie d’une augmentation de taux $t$ correspond à une augmentation globale de  17,6\%.
	Montrer que $t$ est solution de l'équation $\numprint{1.05}(1+t)=\numprint{1.176}$
	et trouver $t$. 
\end{exercise}
\begin{exercise} \hfill
	
	\noindent Une diminution de  15 \%  suivie d’une diminution de taux $t$ correspond à une diminution globale de 32 \%. Trouver $t$.
\end{exercise}
\begin{exercise}[\emoji{cactus}]\hfill
	
	\noindent Après deux augmentations successives de taux $t$, le prix d’un produit a globalement augmenté de 32,25\%. 
	Trouver $t$.
\end{exercise} 
\begin{exercise}[\emoji{cactus}] \hfill 
	
	\noindent	Après une augmentation de  taux  $t$ suivie d’une baisse de taux $t$, le prix d'une chemise a diminué de  4\%.  
	\begin{enumerate}[label={\arabic*)},leftmargin=*] 
		\item Montrer que $\numprint{0.96}=1-t^2$.  
		\item Trouver $t$.
	\end{enumerate} 
\end{exercise}  
\begin{eBox}
	\textbf{Vrai ou Faux ?} \og une augmentation de $t\%$ suivie d'une diminution de $t\%$ est toujours une diminution\fg{}.
\end{eBox} 


\begin{problem}  \hfill
	
	\noindent On aggrandit les côtés d'un carré de 20\%, quel est le taux d'agrandissement des aires.
\end{problem}

\begin{problem} \hfill
	
	\noindent Soit un carré de côté $c$. On augmente deux côtés opposés d'un taux $t_1$, et les deux autres d'un taux $t_2$. 
	\begin{enumerate}[label=\arabic*), leftmargin=*]
		\item Exprimer l'aire du rectangle obtenu en fonction de $c$. 
		\item En déduire que le taux d'évolution de l'aire par cette transformation est $t_1+t_2+t_1t_2$.
		\item Si $t_1=25\%$ et $t_2=-22\%$, l'aire a-t-elle augmenté ou diminué ?
		\item Si $t_1=15\%$. Comment choisir $t_2$ pour que l'aire soit conservée ?
	\end{enumerate}
\end{problem}

\begin{problem}  \hfill
	
	\noindent Deux offres sont proposées pour une bouteille de lessive : 15\% de produit en plus, ou 15\% de réduction sur le prix.	Laquelle choisissez-vous et pourquoi ?
\end{problem}
\begin{problem} \label{chapter11:problem:01}  \hfill
	
	\noindent	Le coût du pétrole représente 24\% du coût de production de l'essence, qui représente lui-même 35\% du prix de l'essence. Si le coût du pétrole augmente de 10\%, et que les autres composants du prix de l'essence ne subissent aucun changement, quel sera le pourcentage d'augmentation du prix de l'essence. 
\end{problem}
\begin{problem}\label{chapter11:problem:02} \hfill
	
	\noindent Après une année d'entrainement, Minnie augmente sa vitesse moyenne au  Marathon de Londres de 25\%. Quel est le taux de diminution de son temps total ?
\end{problem}
\begin{problem}\label{chapter11:problem:03}  \hfill
	
	\noindent Coincée dans le trafic, Emilie's met 25\% de plus que d'habitude pour arriver chez elle. Quel est le pourcentage de réduction de sa vitesse par rapport à un trajet sans embouteillage.
\end{problem}

\begin{problem} \label{chapter11:problem:04} \hfill
	
	\noindent	Si un tapis roulant avançait 10\% plus rapidement, alors le trajet prendrait 5 s de moins. Quelle est la durée du trajet ?
\end{problem}
\begin{proof}[indication pour le problème \ref{chapter11:problem:01}] $x=$ prix de l'essence. Quel est le coût du pétrole ? 
\end{proof} 
\begin{proof}[indication pour le  problème \ref{chapter11:problem:04}]    Traduire l'énoncé en une équation utilisant $l$ et $t$ la longueur du tapis roulant et le temps normal de trajet.
\end{proof}

%----------------------------------------------------------------------------------------
%	 
%----------------------------------------------------------------------------------------

%\newpage 
\begin{exercise}\hfill 
	
\begin{multicols}{2}
\begin{lstlisting}[ language=python , 
%	title=Script C,   
linewidth=0.99\columnwidth,
style=python-code,    
aboveskip=0.0\topskip,
belowskip=0.0\topskip,
xleftmargin=1.2em,
%	lineskip=1.0em,
tabsize=4, 
%numbers=none,
%	firstnumber=10  
] 
def mafonction(t1, t2) :
	cm1 = 1 + t1 / 100
	cm2 = 1 + t2 / 100
	cmg = cm1 * cm2
	tg = (cmg - 1) * 100
	return tg
	
def cm(t) :
	cm = ...
	return ...
def reciproque(t) :
	...
	...
	return ...
\end{lstlisting}
\end{multicols}
\begin{enumerate}[label=\arabic*), leftmargin=*]
	\item Que retourne l'instruction \lstinline|mafonction(100,-29)| ?
	\item Quel est l'objet de la fonction \lstinline|mafonction()| ?
	\item Complétez les lignes 9 et 10 afin que la fonction d'appel \lstinline|cm()| et d'argument \lstinline|t| retourne le coefficient multiplicateur de l'évolution de taux \lstinline|t| donné en \%.
	\item Complétez les lignes 12, 13 et 14 afin que la fonction d'appel \lstinline|reciproque()| et d'argument \lstinline|t| retourne le taux reciproque de l'évolution de taux \lstinline|t| donné en \%.
\end{enumerate}
\end{exercise} 

%----------------------------------------------------------------------------------------
%	cours
%----------------------------------------------------------------------------------------

%\newpage 
%\loadgeometry{widelayout}  
%\pagestyle{scrheadings} % Print headers again  
%\KOMAoptions{
%	headwidth=textwithmarginpar,
%	headsepline=true,
%	headsepline= 0.5pt,
%	footwidth=textwithmarginpar,
%} 
%\onehalfspacing

%----------------------------------------------------------------------------------------
%	 
%----------------------------------------------------------------------------------------
\cleardoublepage 
%----------------------------------------------------------------------------------------
%	
%----------------------------------------------------------------------------------------
\end{document}
